\documentclass[11pt]{article}
\usepackage[margin=1in]{geometry}
\usepackage{amsmath,amssymb}
\usepackage{enumitem}
\usepackage{graphicx}
\usepackage{array}
\usepackage{booktabs}
\usepackage{hyperref}

\title{Wasm to NeoVM Translation Specification}
\author{neo-llvm Project}
\date{\today}

\begin{document}
\maketitle
\tableofcontents
\newpage

\section{Scope}
This specification defines the normative behaviour of the \texttt{wasm-neovm} translator that lowers WebAssembly (Wasm) modules into NeoVM scripts and emits the accompanying NEF and manifest artefacts. It covers:
\begin{itemize}[nosep]
    \item Accepted Wasm module structure, imports/exports, and custom sections.
    \item Translation of supported instructions to NeoVM opcodes.
    \item Runtime helper synthesis for memory, tables, and data/element segments.
    \item Manifest generation, overlay application, safe-flag propagation, and validation.
\item NEF emission with metadata (source URL, method tokens).
\end{itemize}

\section{Definitions}
\begin{description}[style=unboxed,leftmargin=0cm]
    \item[Wasm module] A binary payload conforming to the WebAssembly core spec.
    \item[NeoVM script] A byte sequence of NeoVM opcodes produced by the translator.
    \item[Manifest overlay] A JSON fragment (embedded or external) merged into the generated manifest.
    \item[Translation configuration] The tuple of contract name and optional overlays supplied to the translator.
\end{description}

\section{Pipeline Overview}
\subsection{Stages}
\begin{enumerate}[label=\arabic*., nosep]
    \item \textbf{Wasm Parsing (Frontend)}: Parse sections, record type signatures, imports, exports, memories, tables, globals, data/element segments, and embedded custom sections (\texttt{neo.manifest}, \texttt{neo.methodtokens}, \texttt{neo.source}).
    \item \textbf{Translation}: Lower functions to NeoVM opcodes, synthesise runtime helpers, accumulate method offsets and metadata.
    \item \textbf{Manifest Assembly}: Build a baseline manifest from translated exports, merge overlays, propagate safe flags, and ensure ABI parity.
    \item \textbf{NEF Emission}: Serialize script + metadata into NEF, including compiler ID, optional source URL, and method tokens.
\end{enumerate}

\subsection{Data Flow}
\begin{center}
\begin{tabular}{>{\raggedright}p{3.5cm} >{\raggedright}p{5cm} >{\raggedright\arraybackslash}p{4cm}}
\toprule
Stage & Inputs & Outputs \\
\midrule
Frontend & Wasm bytes, \texttt{ModuleTypes} accumulator & IR: types, imports, exports, tables, globals, segments, custom sections \\
Translation & IR, \texttt{RuntimeHelpers} & Script bytes, translated methods (name, params, return, offset), method tokens, inferred source \\
Manifest & Methods, overlays & JSON manifest (validated), safe flags applied \\
NEF Writer & Script, manifest metadata & \texttt{*.nef}, \texttt{*.manifest.json} \\
\bottomrule
\end{tabular}
\end{center}

\section{Configuration Interfaces}
\subsection{Library Entry Point}
\begin{itemize}[nosep]
    \item \texttt{translate\_with\_config(bytes, TranslationConfig)} is the normative API. It requires:
          \begin{itemize}[nosep]
              \item \texttt{contract\_name}: manifest \texttt{name} field.
              \item Optional \texttt{ManifestOverlay} (value + label) merged during manifest assembly.
          \end{itemize}
    \item \texttt{translate\_module(bytes, name)} is a convenience wrapper equivalent to calling \texttt{translate\_with\_config} with no overlays.
\end{itemize}

\subsection{CLI}
\begin{itemize}[nosep]
    \item \texttt{--input}: Wasm module path (required).
    \item \texttt{--nef}, \texttt{--manifest}: Output paths (default: derived from input basename).
    \item \texttt{--name}: Contract name (defaults to \texttt{Contract}).
    \item \texttt{--manifest-overlay}: Path to an external overlay merged with embedded overlays.
    \item \texttt{--source-url}: Overrides the source URL used in NEF metadata.
    \item \texttt{--compare-manifest}: Optional reference manifest; translation fails (after printing a diff) when the generated manifest diverges from this file.
\end{itemize}

\section{Wasm Module Acceptance Rules}
\subsection{Types and Functions}
\begin{itemize}[nosep]
    \item Function types: parameters/results limited to \texttt{i32} and \texttt{i64}. Multi-value returns are rejected.
    \item Imports: only function imports are permitted; globals, memories, and tables cannot be imported.
    \item Exports: only functions are exported. Duplicate export aliases are allowed; all aliases must reference translated functions.
\end{itemize}

\subsection{Memories}
\begin{itemize}[nosep]
    \item At most one memory; \texttt{memory64} and shared memories are rejected.
    \item Memory index must be zero for all instructions; non-zero indices are rejected.
\end{itemize}

\subsection{Tables}
\begin{itemize}[nosep]
    \item Element type must be \texttt{funcref}; \texttt{externref} and other reference types are rejected.
    \item \texttt{table64} and shared tables are rejected.
    \item Declared element segments must be active or passive; declared-only segments are rejected.
\end{itemize}

\subsection{Globals}
\begin{itemize}[nosep]
    \item Global value types restricted to \texttt{i32} and \texttt{i64}.
    \item Immutable globals may be constant-folded; mutable globals lose constant-propagation after writes.
\end{itemize}

\subsection{Custom Sections}
\begin{itemize}[nosep]
    \item \texttt{neo.manifest}: concatenated JSON fragments merged into the manifest.
    \item \texttt{neo.methodtokens}: concatenated JSON fragments parsed into method tokens and optional source URL.
    \item Unknown sections are ignored.
\end{itemize}

\subsection{Reserved Import Modules}
\begin{itemize}[nosep]
    \item \texttt{syscall}: names must match Neo interop descriptors (for example, \texttt{System.Storage.Get}). Each import lowers to a \texttt{SYSCALL} opcode carrying the hashed descriptor bytes.
    \item \texttt{neo}: exposes DevPack-friendly aliases (\texttt{neo::storage\_read}, \texttt{neo::notify}, \texttt{neo::call\_contract}, ...). The translator resolves these aliases to canonical descriptors before emitting \texttt{SYSCALL}, keeping existing contracts source-compatible.
    \item \texttt{opcode}: imports map directly to NeoVM opcodes. Literal operands are supplied via Wasm immediates; the special helpers \texttt{opcode::RAW} and \texttt{opcode::RAW4} append arbitrary bytes to the script for advanced patterns.
    \item \texttt{env}: limited to C runtime memory utilities (\texttt{memcpy}, \texttt{memmove}, \texttt{memset}). Each function must accept three \texttt{i32} parameters (destination, source-or-value, length) and return \texttt{i32}. The translator synthesises helper stubs that operate on the linear memory buffer so contracts compiled with \texttt{-nostdlib -fno-builtin} remain self-contained.
\end{itemize}

\section{Compatibility Matrix}
This matrix summarises the currently supported Wasm surface. Inputs outside these bounds fail translation with descriptive diagnostics.

\begin{table}[h]
\centering
\begin{tabular}{@{}p{3.6cm}p{9cm}@{}}
\toprule
Area & Status / Notes \\
\midrule
Linear memory & Exactly one 32-bit, non-shared memory is allowed. Passive data segments require a defined memory; multi-memory and shared memories are rejected. \\
Tables and references & Tables must be \texttt{funcref}; \texttt{table64}, shared tables, and reference types beyond \texttt{funcref} are rejected. Runtime helpers cover \texttt{table.get/set/size/grow/fill/copy}, \texttt{table.init}, and \texttt{elem.drop}. \\
Imports & Only functions may be imported. Reserved modules \texttt{syscall}, \texttt{neo}, \texttt{opcode}, and \texttt{env} receive dedicated lowering; importing globals, memories, or tables is unsupported. \\
Globals & Value types must be \texttt{i32}/\texttt{i64}. Initialisers are constant-folded and stored in static slots; mutable globals lose literal tracking once written. \\
Function signatures & Exported ABI methods and \texttt{call\_indirect} signatures accept \texttt{i32}/\texttt{i64} parameters and at most one result. Typed \texttt{select} likewise requires a single \texttt{i32}/\texttt{i64} result. \\
Numeric features & Integer arithmetic, comparisons, conversions, bit counts, and shifts are supported. Floating-point, SIMD, atomics, and threads are not yet implemented and produce explicit errors. \\
\bottomrule
\end{tabular}
\end{table}

\section{Instruction Lowering}
\subsection{Supported Operators}
Each operator lowers to NeoVM opcodes while preserving Wasm stack discipline. Table~\ref{tab:ops} enumerates the supported Wasm instructions together with their NeoVM encoding. Unless stated otherwise, both \texttt{i32} and \texttt{i64} variants are supported and share the same emission logic.

\begin{table}[h]
\centering
\begin{tabular}{@{}lll@{}}
\toprule
Wasm Operator & NeoVM Emission & Notes \\
\midrule
\texttt{i32.const}/\texttt{i64.const} & Smallest \texttt{PUSHINT*} opcode & Literal tracking retained \\
\texttt{local.get/set/tee} & \texttt{LDARG*}/\texttt{STARG*}, \texttt{LDLOC*}/\texttt{STLOC*} & Params use \texttt{LDARGn}, locals use \texttt{LDLOCn} \\
\texttt{global.get/set} & \texttt{LDSFLD*}/\texttt{STSFLD*} & Globals live in static slots \\
\texttt{add, sub, mul} & \texttt{ADD}/\texttt{SUB}/\texttt{MUL} & Shared between 32/64-bit \\
\texttt{and, or, xor} & \texttt{AND}/\texttt{OR}/\texttt{XOR} & Bitwise ops \\
\texttt{shl, shr\_s, shr\_u} & \texttt{SHL}/\texttt{SHR} + masks & Shift counts masked to width, logical shifts zero-extend \\
\texttt{rotl, rotr} & \texttt{ROT} sequences & Implemented via stack shuffles and mask \\
\texttt{eq, ne, lt, gt, le, ge} & \texttt{EQUAL}/\texttt{NOTEQUAL}/\texttt{NUMERIC} compares & Unsigned compares mask both operands \\
\texttt{eqz} & Compare against \texttt{PUSH0} & Result is 0/1 \\
\texttt{div\_s, div\_u, rem\_s, rem\_u} & \texttt{DIV}/\texttt{MOD} + masks & Division by zero traps, unsigned variants zero-extend \\
\texttt{clz, ctz, popcnt} & Helper calls via \texttt{RuntimeHelpers} & Literal folding when both operands are constants \\
\texttt{i32.wrap\_i64}, \texttt{i64.extend\_i32\_{s,u}} & Masking + sign-extension helpers & \texttt{emit\_sign\_extend}/\texttt{emit\_zero\_extend} \\
\texttt{i32.extend8/16\_s}, \texttt{i64.extend8/16/32\_s} & Dedicated sign-extension helpers & Matches Wasm semantics \\
\texttt{select}/typed \texttt{select} & \texttt{JMPIFNOT\_L}/\texttt{DROP}/\texttt{NIP} & At most one result value \\
\texttt{block, loop, if, br, br\_if} & \texttt{JMP\_L}/\texttt{JMPIF\_L} + fix-ups & Preserves stack height guarantees \\
\texttt{br\_table} & Cascading \texttt{JMPIF\_L} with default handler & Selector duplicated, fall-through traps \\
\texttt{call} & \texttt{CALL\_L} + fix-ups & Arguments evaluated left-to-right \\
\texttt{call\_indirect} & Runtime table helper + bounds checks & Null funcref traps \\
\texttt{memory.load/store} & Helper stubs from \texttt{RuntimeHelpers} & Bounds-checked, little-endian \\
\texttt{memory.size}/\texttt{memory.grow} & Helper stubs & Grow returns previous page count or \texttt{-1} \\
\texttt{memory.fill}, \texttt{memory.copy}, \texttt{memory.init}, \texttt{data.drop} & Bulk helper stubs & Active segments applied on first init, passive segments tracked via drop slots \\
\texttt{table.get/set/size/grow/fill/copy}, \texttt{table.init}, \texttt{elem.drop} & Table helper stubs & Tables backed by NeoVM arrays, funcref-only \\
\texttt{ref.null}, \texttt{ref.func}, \texttt{ref.is\_null}, \texttt{ref.eq}, \texttt{ref.as\_non\_null} & Integer sentinels (\texttt{-1} for null) + runtime assertions & Funcref handles only \\
\texttt{drop} & \texttt{DROP} (elided for literals) & Compile-time literals may be removed \\
\texttt{return}/implicit fallthrough & \texttt{RET} & Each function concludes with explicit return \\
\texttt{unreachable} & \texttt{ABORT} & Triggers NeoVM fault \\
\bottomrule
\end{tabular}
\caption{Representative Wasm to NeoVM operator mappings}
\label{tab:ops}
\end{table}

\subsection{Control Flow and Imports}
\begin{itemize}[nosep]
    \item Control frames (function, block, loop, if) track stack height and expected result arity. Branch targets patch the saved frame offsets with relative \texttt{JMP} immediates.
    \item Indirect calls (\texttt{call\_indirect}) first route through a table helper that validates the table index, traps on null funcrefs, and finally dispatches via \texttt{CALL\_L}.
    \item Imports from the reserved modules \texttt{syscall} and \texttt{neo} lower to the NeoVM \texttt{SYSCALL} opcode with the hashed descriptor bytes. The \texttt{opcode} module exposes raw opcodes (plus \texttt{RAW}/\texttt{RAW4} helpers), and the \texttt{env} module supports \texttt{memcpy}/\texttt{memmove}/\texttt{memset} when compiled with \texttt{-nostdlib}.
\end{itemize}

\subsection{Numeric Operations}
\begin{itemize}[nosep]
    \item Arithmetic, logic, and comparison instructions operate on full-width NeoVM integers. Unsigned variants mask both operands using \texttt{emit\_zero\_extend} before performing the signed comparison or division.
    \item Bit-counting operators (\texttt{clz}/\texttt{ctz}/\texttt{popcnt}) fold constant-only inputs at translation time; dynamic operands dispatch to compact helper routines emitted once per module.
    \item Conversion instructions (\texttt{i32.wrap\_i64}, \texttt{i64.extend\_i32\_{s,u}}, and the sign-extension family) map to reusable helper sequences that mask and/or shift the stack value as required.
\end{itemize}

\subsection{Unsupported Features}
\begin{itemize}[nosep]
    \item Floating-point, SIMD, atomics, threads.
    \item Multiple memories, \texttt{memory64}, shared memory, shared tables.
    \item Reference types other than \texttt{funcref}; declared element segments.
    \item Multi-value returns.
\end{itemize}

\section{Runtime Helpers}
\subsection{Memory Helpers}
\begin{itemize}[nosep]
    \item \textbf{Linear memory representation}: the sole Wasm memory is materialised as a NeoVM \texttt{BUFFER} stored in a static slot. Slot 0 holds the buffer, slot 1 the current page count, slot 2 the allocated capacity, and slot 3 the pending initialisation flag.
    \item \textbf{Init}: lazily allocates the buffer on the first memory access, applies active data segments exactly once, and caches the resulting pointer/page metadata.
    \item \textbf{Load/Store}: obtain the buffer from slot~0, bounds-check offset + width against the current length, then perform little-endian loads/stores using \texttt{SUBSTR}, \texttt{SETITEM}, and \texttt{PACK}.
    \item \textbf{Grow}: allocates a new buffer when requested pages exceed the current capacity, copies the existing bytes via \texttt{MEMCPY}, updates page counters, and returns the previous page count or \texttt{-1} if the new size would overflow 32-bit space.
    \item \textbf{Fill/Copy/Init/Drop}: honour passive and active data segments by copying the segment bytes into temporary static slots, using \texttt{MEMCPY} for \texttt{memory.copy}/\texttt{memory.init}, and flipping drop-slot booleans when \texttt{data.drop} is invoked.
\end{itemize}

\subsection{Table Helpers}
\begin{itemize}[nosep]
    \item Tables are represented as NeoVM arrays stored in static slots; funcref indices are encoded as signed integers with \texttt{-1} denoting null.
    \item \textbf{Get/Set/Size/Grow/Fill/Copy}: all operations enforce bounds checking. \texttt{table.grow} respects the declared maximum; new entries are initialised to null.
    \item \textbf{InitFromPassive}/\textbf{ElemDrop}: passive element segments are cached in static slots (value + drop flag). \texttt{table.init} copies them into the target table slice and \texttt{elem.drop} flips the drop flag.
\end{itemize}

\section{Manifest Generation}
\subsection{Baseline Manifest}
The translator synthesizes a manifest with:
\begin{itemize}[nosep]
    \item \texttt{name}: contract name passed via configuration.
    \item \texttt{abi.methods}: one entry per exported function (alias preserved), including parameter types (\texttt{Integer}) and return types (\texttt{Integer}/\texttt{Void}).
    \item \texttt{features}: populated with \texttt{storage} when storage is used.
    \item \texttt{extra}: author/description defaults.
\end{itemize}

\subsection{Overlay Application}
\begin{itemize}[nosep]
    \item Overlays come from embedded \texttt{neo.manifest} sections and optional external overlays (\texttt{ManifestOverlay}).
    \item Merging is deep and deduplicates \texttt{supportedstandards}, \texttt{permissions}, \texttt{trusts}, and method entries by name.
    \item Safe flags are propagated across duplicate method entries.
    \item ABI parity check: if overlays introduce or remove methods relative to translated exports, translation fails with the overlay label in the diagnostic.
    \item DevPack macros mirror the JSON structure: \texttt{neo\_event} emits canonical \texttt{abi.events} entries (Boolean/Integer/ByteArray/\dots), while \texttt{neo\_permission!}, \texttt{neo\_supported\_standards!}, and \texttt{neo\_trusts!} append the corresponding manifest fragments before merging.
    \item The translator infers `features.storage` automatically whenever translated code invokes a `System.Storage.*` syscall, and enables `features.payable` when an exported entry point named `onPayment`, `onNEP17Payment`, or `onNEP11Payment` is present, so overlays rarely need to toggle either flag manually.
\end{itemize}

\section{NEF Emission}
\subsection{Metadata}
\begin{itemize}[nosep]
    \item Compiler ID: fixed string embedded in NEF header.
    \item Source URL: taken from \texttt{neo.source} custom section or CLI/config override (\texttt{TranslationConfig}); optional.
    \item Method tokens: parsed from custom sections, inferred from contract calls, and deduplicated before writing.
\end{itemize}

\subsection{File Format Requirements}
\begin{itemize}[nosep]
    \item NEF header includes magic, compiler, source URL, reserved byte, method tokens, script, checksum.
    \item Manifest and NEF must remain in sync; the translator writes both or fails.
\end{itemize}

\section{Errors and Diagnostics}
Translation aborts with descriptive errors on:
\begin{itemize}[nosep]
    \item Unsupported Wasm features (floating point, multi-memory, non-funcref tables, etc.).
    \item ABI mismatch introduced by overlays.
    \item Out-of-bounds segment initialisation or table/memory access during helper generation.
    \item Invalid custom-section payloads for \texttt{neo.manifest} or \texttt{neo.methodtokens}.
\end{itemize}

\section{Conformance}
An implementation conforms to this specification if, given any Wasm module satisfying the acceptance rules, it:
\begin{itemize}[nosep]
    \item Produces a NeoVM script whose observable behaviour matches the specified lowering.
    \item Emits NEF and manifest artefacts that reflect the translated ABI and overlays, with safe flags and metadata propagated.
    \item Rejects inputs that violate acceptance rules or manifest parity, emitting diagnostics that identify the violating feature or overlay.
\end{itemize}

\end{document}
